%% toc: P15 164 164.12
\chapter{Confirmation and Statistical Estimation: Two Source Comparisons}
\label{sec:p15-w02b-p15-025-comparison}

This chapter makes two separately attributed comparisons. The first,
on confirmation, sets two quantitative notions of how data bear on a
theory side by side. Hempel and Oppenheim, in Part~IV of their 1948
paper, measure the \emph{systematic power} of a theory with respect to
data and show that, for every range measure, it is the dual of logical
probability \cite[pp.~164--173]{HempelOppenheim1948Explanation}.
Merin's chapter on evidential relevance measures how far learning a
\emph{source} proposition changes the credence of a \emph{target}
proposition, and separates such changes from the posterior credence
itself \cite[pp.~265--275]{Merin2021Relevance}. The two sources are
presented in full in Chapters~\ref{sec:p15-systematic-power}
and~\ref{sec:p15-rationality-4-3}, with Merin's relevance axioms and
Carnap's covariance measure continued in
Chapter~\ref{sec:p15-w02b-p15-024-relevance}. This chapter supplies only
what neither presentation contains: the exact formal relations between
the two families of quantities, the conditions under which they agree,
and the comparisons on which they come apart. The second comparison,
from \cref{sec:p15-w02b-p15-025-carnap98} on, takes Carnap's 1950
account of how confirmation-based estimation stands to the statistical
methods of Fisher, Neyman and Pearson, and Wald, and gives each of his
formal contrasts an exact model.

\paragraph{Inputs and conclusions.}
The input is a fixed language in the sense of
\cref{def:p15-power-language}, a range measure \(r\) on its sentences
(\cref{def:p15-power-range-measure}), a theory \(T\) and a data sentence
\(S\). The conclusions are identities and equivalences that hold for
every such range measure, together with explicit finite models that
separate the notions. Each identity is a mathematical consequence of the
stated definitions. Which range measure, and therefore which numbers, a
user should adopt is a separate question that both sources leave open
(\cref{sec:p15-w02b-p15-025-limits}).

\section{What each source measures}
\label{sec:p15-w02b-p15-025-sources}

\paragraph{Hempel and Oppenheim.}
Their starting point is that laws and theories connect data so that
some data can be derived from others; they treat explanation and
prediction alike as deductive systematization, the difference being only
whether the derived data were already known when the derivation was
made, and they adopt the neutral name \emph{systematic power} for the
measure they seek \cite[p.~164]{HempelOppenheim1948Explanation}. The
first definition, (8.1), counts the minimal sentences of the data's
content that can be derived from the rest by means of the theory
\cite[pp.~166--167]{HempelOppenheim1948Explanation}. Section~9 then
calls logical probability, or degree of confirmation, the central concept
of inductive logic, writes it as \(p(T,S)=r(T\cdot S)/r(S)\) for a range
measure \(r\), and defines the generalized systematic power
\(s(T,S)=g(T\vee S)/g(S)\) with \(g=1-r\)
\cite[pp.~168--172, (9.3), (9.6), (9.9), (9.10)]{HempelOppenheim1948Explanation}.
Their identity (9.11) states \(s(T,S)=p(\neg T,\neg S)\), and their
Dualism Theorem (9.13) transfers every general equality or inequality
about \(p\) to one about \(s\) \cite[pp.~172--173]{HempelOppenheim1948Explanation}.
In note~35 they describe their \(s\) as a quantitative counterpart of
the prediction criterion of confirmation that Hempel had stated in
non-numerical terms \cite[p.~168, n.~35]{HempelOppenheim1948Explanation};
that criterion is \cref{def:p15-hempel-prediction}.

In the notation of \cref{def:p15-power-dual-measures}, these are
\(p_r(T,S)\) and \(s_r(T,S)\). Hempel and Oppenheim use ``degree of
confirmation'' for the conditional quantity \(p\); they do not define a
change in probability produced by the data.

\paragraph{Merin.}
Merin takes relevance to be a change of credence under conditioning:
\(A\) is positive to \(B\) when \(P(B\mid A)>P(B)\), equivalently when
\(P(AB)>P(A)P(B)\) (his Definitions~2 and~3), and a proposition of
probability zero or one is irrelevant to everything (his Theorem~3)
\cite[p.~266]{Merin2021Relevance}. He warns that ``confirmation'',
``corroboration'' and ``support'' need the prefix ``incremental'' if
they are not to be read as the conditional probability \(P(H\mid E)\),
and introduces a directed notation \(R_T(A\Yright B)_P\) with \(A\) in
source position and \(B\) in target position
\cite[p.~268, Definition~6]{Merin2021Relevance}. A relevance measure
must be a function of the joint table of the two propositions and must
represent the sign of relevance by a fixed neutral value
\cite[p.~269, Definition~8]{Merin2021Relevance}. The difference measure
\(R_D\), the ratio \(R_Q\) and the likelihood ratio \(R_L\), with their
transforms, are given in \cref{def:p15-r43-measures}; the axiom POP-D,
which orders sources for a fixed target by their posteriors, is in
\cref{sec:p15-r43-axioms} \cite[pp.~269--271]{Merin2021Relevance}.

\paragraph{The comparison in one sentence.}
Hempel and Oppenheim's \(p\) is what Merin calls a posterior, not an
amount of relevance; their \(s\) is a posterior of a different event,
the conditional probability that the theory fails given that the data
fail. Neither is a relevance measure in Merin's sense. The results below
make this exact and identify which relevance quantity each of them
contains.

\section{One probability model for both sources}
\label{sec:p15-w02b-p15-025-model}

\begin{hbconvention}{def:p15-w02b-p15-025-common-model}{Reading a range measure as a credal probability}
Fix a range measure \(r\) satisfying
\cref{def:p15-power-range-measure}. Let \(\mathcal B\) be the algebra of
sentences of the fixed language taken up to logical equivalence, with
\(\wedge,\vee,\neg\) as meet, join and complement. Put
\[
 P_r([U])=r(U)
\]
for the class \([U]\) of each sentence \(U\). All relevance relations
and measures of Merin's chapter are then evaluated under \(P_r\); for
instance \(\Delta_r(A,B)=r(A\wedge B)-r(A)\,r(B)\), and
\(R_D(S\Yright T)_r=p_r(T,S)-r(T)\) when \(r(S)>0\). Write
\(O(x)=x/(1-x)\) for \(0\le x<1\).

The reading is formal. Hempel and Oppenheim's \(r\) is a measure fixed
on a language before any data are assessed; Merin's \(P\) is the credal
state of a reasoner. Treating the one as the other is a decision of the
user, not a claim of either source.
\end{hbconvention}

\begin{hbtheorem}{thm:p15-w02b-p15-025-common-model}{proved=annex:p15-w02b-p15-025-common-model}{The range measure is a finitely additive probability}
\(P_r\) is well defined on \(\mathcal B\), takes the value \(1\) at the
class of valid sentences and is finitely additive on disjoint classes.
For every sentence \(T\),
\[
 \begin{gathered}
 p_r(T,S)=P_r(T\mid S)\quad\text{when }r(S)>0,\\
 s_r(T,S)=P_r(\neg T\mid\neg S)\quad\text{when }r(S)<1,
 \end{gathered}
\]
and the two conditions on \(r(S)\) are exactly the domains on which
\(p_r\) and \(s_r\) are defined.
\end{hbtheorem}

\section{Systematic power and the prediction criterion}
\label{sec:p15-w02b-p15-025-prediction}

Hempel and Oppenheim's note~35 is a claim about a qualitative criterion.
The following theorem states how far it holds for their original measure
\(s_0\) of \cref{def:p15-power-optimal-split}.

\begin{hbtheorem}{thm:p15-w02b-p15-025-prediction}{proved=annex:p15-w02b-p15-025-prediction}{The original measure is positive exactly when the prediction criterion holds for the content}
Let \(T\) be any sentence and \(K\) a finite set of singular sentences
whose conjunction \(S\) is not valid. Call a report \(E\) \emph{predictive
for} \(T\) when it has a partition \(E=E_1\mathbin{\dot\cup}E_2\) with
\(E_2\neq\varnothing\) such that every member of \(E_2\) follows from
\(E_1\cup\{T\}\) and not from \(E_1\) alone; this is the confirming clause
of \cref{def:p15-hempel-prediction}.
\begin{enumerate}
\item The report \(\mathcal C(S)\) of minimal sentences in the content of
\(S\) is predictive for \(T\) if and only if \(s_0(T,K)>0\).
\item If the report \(K\) itself is predictive for \(T\), then
\(s_0(T,K)>0\).
\item The converse of item~2 fails: for the theory and the data set
\(K_2\) of \cref{thm:p15-power-source-examples}, \(s_0(T,K_2)=25/36\),
while no sentence of \(K_2\) is derivable from the other members together
with \(T\), so \(K_2\) is not predictive for \(T\).
\end{enumerate}
\end{hbtheorem}

Item~1 makes the note's ``quantitative counterpart'' exact: positivity
of \(s_0\) is the prediction criterion applied to the content of the data
rather than to the sentences in which a report happens to be written.
Item~3 shows that this choice matters. The criterion read on the raw
sentences depends on how the data are grouped, which is the defect that
led Hempel and Oppenheim to minimal sentences in the first place
\cite[p.~165]{HempelOppenheim1948Explanation}. The other defects of the
prediction criterion that Hempel identified, including its failure to
pass confirmation to weaker hypotheses, are treated in
\cref{thm:p15-hempel-prediction-gap}; item~1 does not remove them.

\section{The increment hidden in systematic power}
\label{sec:p15-w02b-p15-025-increment}

\begin{hbtheorem}{thm:p15-w02b-p15-025-increment}{proved=annex:p15-w02b-p15-025-increment}{Systematic power above content is a scaled difference measure}
Let \(0<r(S)<1\). Then, for every sentence \(T\),
\[
 s_r(T,S)-g(T)
   =\frac{\Delta_r(S,T)}{r(\neg S)}
   =O\bigl(r(S)\bigr)\,R_D(S\Yright T)_r
   =R_D(\neg S\Yright\neg T)_r .
\]
Consequently:
\begin{enumerate}
\item \(s_r(T,S)>g(T)\), \(=g(T)\) or \(<g(T)\) exactly when \(S\) is
positive, irrelevant or negative to \(T\) under \(P_r\).
\item The quantity \(s_r(T,S)-s_r(T,\bot)\) is the image of
\(p_r(T,S)-p_r(T,\top)\) under the dual transformation of
\cref{thm:p15-power-dualism}, and \(s_r(T,\bot)=g(T)\). Both
differences have the sign of \(\Delta_r(S,T)\) on the domain
\(0<r(S)<1\).
\item If \(r(S)=0\), then \(s_r(T,S)=g(T)\) for every \(T\).
\item The function \(s_r\) itself does not satisfy Merin's Quality
Representation condition: no single value \(\iota\) separates
positive from negative relevance for all theories. With the uniform
measure on three ground atoms \(A,B,C\), \(s_r(A,B)=1/2\) and
\(s_r(A\wedge C,B)=3/4\), although \(B\) is irrelevant to both
targets.
\end{enumerate}
\end{hbtheorem}

The theorem locates Merin's sign of relevance inside Hempel and
Oppenheim's quantity: what exceeds the theory's unconditional content
\(g(T)\) is a difference measure, read with both propositions negated
and rescaled by the prior odds of the data. Merin's polarity laws
(\cref{thm:p15-r43-polarity}) already say that negating both
arguments preserves the sign; item~1 is that law applied to \(s_r\).
Item~4 is why \(s_r\) itself cannot serve as a relevance measure: its
neutral point \(g(T)\) moves with the target.

\section{Maximal power is conclusive relevance of the theory to the data}
\label{sec:p15-w02b-p15-025-conclusive}

\begin{hbtheorem}{thm:p15-w02b-p15-025-conclusive}{proved=annex:p15-w02b-p15-025-conclusive}{When systematic power equals one}
Let \(r(T)>0\) and \(0<r(S)<1\). The following are equivalent:
\begin{enumerate}
\item \(s_r(T,S)=1\);
\item \(r(T\wedge\neg S)=0\);
\item \(T\) is extremely positive to \(S\) under \(P_r\), in the sense
of \cref{def:p15-r43-extreme};
\item \(\neg S\) is extremely positive to \(\neg T\) under \(P_r\).
\end{enumerate}
They hold in particular when \(T\entails S\). Whenever they hold,
\(0<r(T)\le r(S)<1\) and
\[
 \begin{aligned}
 p_r(T,S)&=\frac{r(T)}{r(S)}, &
 R_D(S\Yright T)_r&=\frac{r(T)\,(1-r(S))}{r(S)},\\
 R_Q(S\Yright T)_r&=\frac1{r(S)}, &
 R_L(S\Yright T)_r&=\frac{1-r(T)}{r(S)-r(T)},\\
 s_r(T,S)-g(T)&=r(T),
 \end{aligned}
\]
with \(R_L=+\infty\) when \(r(S)=r(T)\). For a fixed \(S\), the values
of \(p_r\), \(R_D\) and \(R_L\) are strictly increasing in \(r(T)\),
while \(s_r\) and \(R_Q\) take the same value for every theory
satisfying the four conditions.
\end{hbtheorem}

The direction matters. Maximal systematic power is conclusive relevance
\emph{of the theory to the data}, the direction Merin describes when he
notes that reasoning often runs from hypotheses to possible data
\cite[p.~268]{Merin2021Relevance}. The likelihood-ratio measure and
\(R_Z\) take their maximal values exactly at conclusive positive
relevance \cite[p.~272]{Merin2021Relevance}, so \(s_r(T,S)=1\) is the
same event as \(R_L(T\Yright S)_r=+\infty\) under these
nondegeneracy conditions. The relevance of the data to the theory is a
different quantity: for theories that all entail the data it is decided
by the theories' prior range, which systematic power ignores. Merin's
hypothetico-deductive argument in his Section~4, that a hypothesis
entailing a possible outcome is raised by observing it, is the sign of
the displayed value of \(R_D(S\Yright T)_r\)
\cite[pp.~267--268]{Merin2021Relevance}.

\section{Ranking data for a fixed theory}
\label{sec:p15-w02b-p15-025-ranking}

Merin's Theorem~9 says that all measures validating POP-D agree when
several sources are ranked for one target
\cite[p.~271]{Merin2021Relevance}; in this handbook this is
\cref{thm:p15-r43-rankings}. Systematic power does not belong to that
family.

\begin{hbtheorem}{thm:p15-w02b-p15-025-not-popd}{proved=annex:p15-w02b-p15-025-not-popd}{Systematic power reverses a POP-D ranking}
On the eight conjunctions of \(H,E_1,E_2\) and their negations, let
\(P(H)=1/2\), let \(E_1,E_2\) be independent conditionally on \(H\) and
on \(\neg H\), and let
\[
 \begin{gathered}
 P(E_1\mid H)=\tfrac34,\qquad P(E_1\mid\neg H)=\tfrac14,\\
 P(E_2\mid H)=\tfrac7{16},\qquad P(E_2\mid\neg H)=\tfrac1{16}.
 \end{gathered}
\]
Take \(r\) to be this probability on the singular sentences of the
three atoms. Then
\[
 \begin{array}{c|cccc}
  &r(E_i)&p_r(H,E_i)&s_r(H,E_i)&s_r(H,E_i)-g(H)\\\hline
 E_1&1/2&3/4&3/4&1/4\\
 E_2&1/4&7/8&5/8&1/8
 \end{array}
\]
Every measure validating POP-D ranks \(E_2\) strictly above \(E_1\) as
relevance to \(H\); systematic power, and its excess over \(g(H)\), rank
\(E_1\) strictly above \(E_2\). Neither \(s_r(H,\cdot)\) nor
\(s_r(H,\cdot)-g(H)\) is ordinally equivalent, on fixed-target source
rankings, to any measure satisfying POP-D.
\end{hbtheorem}

The reversal has a plain reading. For a fixed theory, systematic power
rewards data whose failure would count heavily against the theory, and
by the identity of \cref{thm:p15-w02b-p15-025-increment} it weights the
posterior gain by the prior odds of the data. Probable data that
modestly raise the theory can therefore outrank improbable data that
raise it more. This is a difference in what is being measured, not an
error in either source: Hempel and Oppenheim do not offer \(s\) as a
measure of incremental support.

\section{Data that are already certain}
\label{sec:p15-w02b-p15-025-certain}

\begin{hbtheorem}{thm:p15-w02b-p15-025-certain}{proved=annex:p15-w02b-p15-025-certain}{Certain data under a logical and under a credal measure}
\begin{enumerate}
\item \(s_r(T,S)\) is defined if and only if \(r(S)<1\).
\item Let the language be singular with ground atoms \(A_1,\ldots,A_n\),
and let \(r\) give every ground assignment positive weight. For singular
\(S\), \(r(S)<1\) if and only if \(S\) is not valid, and
\(0<r(S)<1\) if and only if \(S\) is neither valid nor contradictory.
\item Let \(r(S)>0\) and let \(r_S(U)=r(U\wedge S)/r(S)\) be the range
measure conditioned on \(S\). Then \(r_S\) satisfies
\cref{def:p15-power-range-measure}, \(r_S(S)=1\), \(s_{r_S}(T,S)\) is
undefined for every \(T\), and \(S\) is irrelevant to every sentence
under \(P_{r_S}\).
\end{enumerate}
\end{hbtheorem}

Hempel and Oppenheim treat data that are already known and data not yet
observed alike: explanation and prediction differ for them only in the
time of the derivation \cite[p.~164]{HempelOppenheim1948Explanation}.
Their uniform measure \(r_1\) of (9.4) gives every ground assignment the
same weight \cite[p.~170]{HempelOppenheim1948Explanation}, so by item~2
the original measure is defined for all nonvalid data whether or not
anyone knows them. Item~3 is the situation Merin describes for old
evidence: once a credal state has absorbed the data, the data are
irrelevant to every hypothesis (his Theorem~3), and the reconstructions
he lists work by changing the credal state, the update history or the
measure \cite[pp.~266, 268]{Merin2021Relevance}; see
\cref{sec:p15-r43-old-evidence}. A logical range measure avoids that
predicament only because it is not anybody's current credence. That is
a reason to keep the two readings of \(r\) separate, not a solution of
the old-evidence problem.

\section{Worked example}
\label{sec:p15-w02b-p15-025-example}

\begin{hbtheorem}{thm:p15-w02b-p15-025-example}{proved=annex:p15-w02b-p15-025-example}{Hempel and Oppenheim's illustration under the uniform measure}
Use the language with one constant \(a\) and unary predicates \(P,Q\),
and the uniform measure \(r_1\) on its four ground assignments.
\begin{enumerate}
\item For \(T=\forall x(P(x)\to Q(x))\) and \(K=\{P(a),Q(a)\}\), with
\(S=P(a)\wedge Q(a)\) and the singular projection
\(T^\flat=\neg P(a)\vee Q(a)\) of \cref{thm:p15-power-uniform-formula},
\[
 \begin{gathered}
 s_0(T,K)=s_{r_1}(T^\flat,S)=\tfrac13,\qquad
 g_1(T^\flat)=\tfrac14,\qquad
 p_{r_1}(T^\flat,S)=1,\\
 R_D(S\Yright T^\flat)_{r_1}=\tfrac14,\qquad
 R_L(S\Yright T^\flat)_{r_1}=+\infty .
 \end{gathered}
\]
The data are conclusively positive to \(T^\flat\), while the systematic
power is \(1/3\), and \(1/3-1/4=O(1/4)\cdot1/4\).
\item For \(S'=P(a)\vee Q(a)\), the theories \(T_1=P(a)\wedge Q(a)\) and
\(T_2=P(a)\) both entail \(S'\), and
\[
 \begin{array}{c|cccccc}
  &s_{r_1}&p_{r_1}&R_D&R_Q&R_L&s_{r_1}-g_1\\\hline
 T_1&1&1/3&1/12&4/3&3/2&1/4\\
 T_2&1&2/3&1/6&4/3&2&1/2
 \end{array}
\]
where each relevance value is that of \(S'\) to the theory.
\end{enumerate}
\end{hbtheorem}

The first item is Hempel and Oppenheim's own illustration, in which they
compute \(s(T,K)=1/3\) \cite[p.~167]{HempelOppenheim1948Explanation}.
The relevance values use the projection \(T^\flat\), not \(T\): the
uniform measure \(r_1\) is defined on singular sentences only, and a
range measure on the quantified \(T\) would have to be chosen
separately. The two numbers answer different questions. The data settle
the singular content of the law at the instance \(a\), which is what
\(p\) and \(R_L\) register; they cover only one of the three minimal
pieces of their own content by means of the law, which is what \(s_0\)
registers. The second item shows two theories that systematize the same
data completely and receive the same ratio relevance, but differ in
posterior, in difference relevance and in likelihood ratio.

\section{Limits of the comparison}
\label{sec:p15-w02b-p15-025-limits}

\begin{hbwarning}{def:p15-w02b-p15-025-limits}{What the identities do not settle}
\begin{enumerate}
\item The identities hold for one range measure at a time. Hempel and
Oppenheim end Part~IV by naming the choice of an adequate range measure
among infinitely many formal possibilities, and the specification of
qualitative primitives, as open problems
\cite[p.~173]{HempelOppenheim1948Explanation}. Nothing here selects
\(r\).
\item Carnap's logical probability is extended to sentences that are not
both singular by limiting processes, and the degree of confirmation of
Helmer, Hempel and Oppenheim depends on a statement of empirical
information and, by the authors' own account in note~40, is not a
probability in the strict sense
\cite[pp.~170--171, nn.~38, 40]{HempelOppenheim1948Explanation}. The
results of this chapter assume a finitely additive range measure on the
whole language, as in \cref{def:p15-power-range-measure}, and do not
apply to those constructions without that assumption.
\item Merin's measures are defined for a credal probability and his
axioms are stated over all probability models
\cite[pp.~269--271]{Merin2021Relevance}. A single range measure is one
such model. The comparisons in
\cref{thm:p15-w02b-p15-025-not-popd,thm:p15-w02b-p15-025-increment}
exhibit models; they are not claims about which model is correct.
\item The relevance values computed for a quantified theory use its
singular projection. For a quantified \(T\), relevance under a measure
that also weighs non-singular sentences can differ.
\item No claim is made about how scientists or other reasoners assess
theories. Merin's psychological applications are treated, with their
own limits, in \cref{sec:p15-r43-use}.
\end{enumerate}
\end{hbwarning}

\section{Carnap on confirmation and statistical estimation}
\label{sec:p15-w02b-p15-025-carnap98}

The second comparison concerns statistical inference. Section~98 of
Carnap's \emph{Logical Foundations of Probability} sets the estimation
methods of Fisher, of Neyman and Pearson, and of Wald against an
estimate defined from a degree of confirmation
\cite[sec.~98, pp.~511--520]{Carnap1950Probability}. The
\(c\)-mean estimate itself and Carnap's summary of the three rivals are
in \cref{sec:p15-carnap-general-estimation}; this chapter states the
claims of Section~98 by type and then supplies, for each rival, an exact
formal contrast that Section~98 announces but does not carry out.

\paragraph{Inputs and conclusions.}
The inputs are a statistical model (a family of sampling laws indexed by
a parameter) and, on Carnap's side, a confirmation function that also
assigns values to the parameter hypotheses. The conclusions are
identities and finite models that show where the two kinds of output
coincide, where they differ, and by how much. None of them decides
which method is correct; Carnap himself makes that a question of the
values the methods deliver, to be compared case by case
\cite[sec.~98, pp.~519--520]{Carnap1950Probability}.

\paragraph{Carnap's claims, by type.}
\begin{enumerate}
\item \emph{Formal.} An estimate is to be a mean of the possible values
of a magnitude weighted by their degrees of confirmation
\cite[ch.~IX, p.~511]{Carnap1950Probability}
(\cref{def:p15-carnap-c-mean}). Fisher's maximum likelihood estimate is
the parameter value under which the observed sample has the largest
probability, where that probability is the statistical one, equal to the
degree of confirmation of the sample description given the parameter
statement \cite[sec.~98, p.~515]{Carnap1950Probability}. Neyman's
confidence coefficient is a probability, in the statistical sense, of
correct interval statements in repeated sampling; it is equal to the
confirmation of the statement that the interval contains the parameter
\emph{relative to a hypothesis about the parameter}, a direct
inference, and is not the confirmation of that statement relative to
the observed sample \cite[sec.~98, p.~517]{Carnap1950Probability}.
Wald's minimax rule chooses the decision function whose largest risk
over parameter values is smallest, whereas Carnap's own decision rules
(his Sections~50E and~51A) minimize the risk weighted by degree of
confirmation \cite[sec.~98, pp.~517--518]{Carnap1950Probability}.
\item \emph{Historical.} The estimate of a population variance from a
sample was first the sample variance; Gauss and Bessel proposed a
corrected value with the factor \(n/(n-1)\); most statisticians
followed them, while Fisher's method yields the uncorrected value
\cite[sec.~98, p.~513]{Carnap1950Probability}. The motive for methods
independent of degree of confirmation was, in Carnap's reading, chiefly
negative: rejection of the principle of indifference, which enters
parameter estimation through the Bayes--Laplace theorem
\cite[sec.~98, p.~518]{Carnap1950Probability}.
\item \emph{Normative.} A single choice of confirmation function would
supply one general method for all estimation problems, where current
statistics answers each problem by separate plausibility arguments
\cite[sec.~98, p.~514]{Carnap1950Probability}; and if no adequate
quantitative confirmation concept exists, the independent methods are
presumably the best instruments available
\cite[sec.~98, p.~518]{Carnap1950Probability}.
\item \emph{Deferred.} Carnap says that serious reasons for doubting
maximum likelihood values on small samples, and his comparison of
\(c^*\) and Wald's rule with other methods, will be given in a second
volume \cite[sec.~98, pp.~515, 518--520]{Carnap1950Probability}. He
also restricts the available theory to the simple languages of the
book, leaving the full quantitative language of physics as a future
task \cite[sec.~98, p.~519]{Carnap1950Probability}.
\end{enumerate}
The results below give exact content to the formal claims, in models
where each of them is meaningful. They are not taken from Carnap's
text, which supplies no such calculations in Section~98.

\section{Maximum likelihood and the \texorpdfstring{\(c^*\)}{c*}-estimate of a frequency}
\label{sec:p15-w02b-p15-025-ml}

Carnap's announced comparison concerns chiefly the estimate of the
relative frequency of a property from an observed sample
\cite[sec.~98, p.~520]{Carnap1950Probability}. On his side the value is
the singular predictive confirmation of \cref{thm:p15-carnap-predictive},
which is also the \(c^*\)-mean estimate of the relative frequency in an
unobserved list. On Fisher's side the model is the independent
Bernoulli model of \cref{def:p15-fe-binomial}, whose maximum likelihood
estimate is the observed frequency (\cref{thm:p15-fe-binomial}). To
compare the two outputs one must read Fisher's estimate of the success
probability as an estimate of the frequency in further trials. That
identification is a modelling choice, made here explicitly.

\begin{hbtheorem}{thm:p15-w02b-p15-025-ml-cstar}{proved=annex:p15-w02b-p15-025-ml-cstar}{Distance between the two frequency estimates}
In the setting of \cref{thm:p15-carnap-predictive}, let \(M\) have width
\(w\) in a division of total width \(k\), with \(0<w<k\), and let the
evidence report \(s_M\) individuals with \(M\) among \(s\) observed
individuals. Write \(c^*=(s_M+w)/(s+k)\).
\begin{enumerate}
\item If \(s\ge1\), then
\[
 c^*-\frac{s_M}{s}=\frac{ws-ks_M}{s(s+k)},\qquad
 \Bigl|c^*-\frac{s_M}{s}\Bigr|\le\frac{\max(w,k-w)}{s+k},
\]
and the difference is zero exactly when \(s_M/s=w/k\). The bound is
attained at \(s_M=0\) when \(w\ge k-w\) and at \(s_M=s\) when
\(w\le k-w\).
\item If \(s\ge1\) and \(s_M=0\), the maximum likelihood estimate of the
Bernoulli parameter is \(0\), while \(c^*=w/(s+k)>0\). If \(s_M=s\),
it is \(1\), while \(c^*=(s+w)/(s+k)<1\).
\item If \(s=0\), every parameter value maximizes the Bernoulli
likelihood, so there is no unique maximum likelihood estimate, while
\(c^*=w/k\).
\end{enumerate}
\end{hbtheorem}

The distance is at most of order \(1/s\), which is an exact form of
Carnap's remark that maximum likelihood values are often acceptable
approximations for large samples \cite[sec.~98, p.~515]{Carnap1950Probability}.
The disagreement is greatest at the extreme counts, where the observed
frequency gives the next case probability zero or one.

\begin{hbexample}{def:p15-w02b-p15-025-small-sample}{Four trials without a success}
Take one primitive property, so \(k=2\) and \(w=1\), and four observed
individuals none of which has \(M\). The maximum likelihood estimate is
\(0\); \(c^*=1/6\), which attains the bound of
\cref{thm:p15-w02b-p15-025-ml-cstar}. For the same count,
\cref{thm:p15-confidence-credibility-boundary} computes the posterior
content \(1-20^{-5/4}\) of the exact \(9/10\) confidence interval under
the uniform prior, which is not \(9/10\). Three outputs from the same
data answer three questions: which parameter makes the data most
probable, how far \(M\) at the next individual is confirmed, and which
interval carries a repeated-sampling guarantee.
\end{hbexample}

\section{The Gauss--Bessel correction}
\label{sec:p15-w02b-p15-025-variance}

Carnap uses the estimate of a variance as his example of a controversy
settled by plausibility rather than proof
\cite[sec.~98, pp.~513--514]{Carnap1950Probability}. The next result
shows that each competing value is optimal under a different stated
criterion, so that none of these properties settles the choice alone.

\begin{hbtheorem}{thm:p15-w02b-p15-025-variance}{proved=annex:p15-w02b-p15-025-variance}{Three criteria, three multipliers}
Let \(n\ge2\) and \(\nu=n-1\), and let \(\bar X\) and
\(V=\sum_{i=1}^n(X_i-\bar X)^2\) be as in \cref{def:p15-student-model}.
\begin{enumerate}
\item If \(X_1,\ldots,X_n\) are independent with common mean \(\mu\)
and common finite variance \(\sigma^2\), then
\(\mathbb E[V/n]=\nu\sigma^2/n\) and \(\mathbb E[V/\nu]=\sigma^2\).
So the corrected value \(V/(n-1)\) is unbiased and the sample variance
\(V/n\) is not.
\item In the normal model of \cref{def:p15-student-model}, if the
observed values are not all equal, the likelihood has the unique
maximizer \((\mu,\sigma^2)=(\bar x,v/n)\), where \(v\) is the observed
value of \(V\). If they are all equal, the likelihood is unbounded and
has no maximizer.
\item In the same normal model, for every constant \(a>0\),
\[
 \mathbb E_{\mu,\sigma}\bigl[(aV-\sigma^2)^2\bigr]
 =\sigma^4\bigl(a^2\nu(\nu+2)-2a\nu+1\bigr),
\]
which is minimized uniquely at \(a=1/(n+1)\). The values at
\(a=1/(n+1),1/n,1/(n-1)\) are
\[
 \frac{2\sigma^4}{n+1}<\frac{(2n-1)\sigma^4}{n^2}<\frac{2\sigma^4}{n-1}.
\]
\end{enumerate}
\end{hbtheorem}

Item~2 confirms Carnap's attribution that Fisher's method gives the
uncorrected value \cite[sec.~98, p.~513]{Carnap1950Probability}.
Unbiasedness selects \(V/(n-1)\), likelihood selects \(V/n\), and mean
squared error among multiples of \(V\) selects \(V/(n+1)\), which
neither side of Carnap's historical controversy proposed. Items~1
and~3 are statements about repeated sampling at a fixed parameter; a
\(c\)-mean estimate of \(\sigma^2\) would need a confirmation function
on the parameters and is not computed here.

\section{Confidence coefficient and degree of confirmation}
\label{sec:p15-w02b-p15-025-confidence}

\begin{hbdefinition}{def:p15-w02b-p15-025-joint}{A sampling model joined to a confirmation function}
Let \(\Theta\) and \(\mathsf X\) be finite sets, \(P_\theta\) a
probability on \(\mathsf X\) for each \(\theta\in\Theta\), and
\(C:\mathsf X\to2^\Theta\) an unrandomized confidence procedure in the
sense of \cref{def:p15-confidence-procedure}, with target
\(\tau(\theta)=\theta\). Its coverage at \(\theta\) is
\[
 \gamma_C(\theta)=P_\theta\{x:\theta\in C(x)\}.
\]
A confirmation function that gives the parameter hypotheses the values
\(\pi(\theta)\), and gives a sample description, relative to the
hypothesis that the parameter is \(\theta\), the value \(P_\theta(x)\),
determines the joint probability
\[
 Q_\pi(\theta,x)=\pi(\theta)P_\theta(x)
\]
on \(\Theta\times\mathsf X\). Write \(Q_\pi(x)=\sum_\theta Q_\pi(\theta,x)\).
\end{hbdefinition}

The second clause is Carnap's requirement that the statistical
probability given a parameter value equal the confirmation in the
corresponding direct inference \cite[sec.~98, p.~517]{Carnap1950Probability}.
The values \(\pi\) are not supplied by any sampling model; they are the
extra input that Carnap's method needs and that the methods he
discusses decline to use.

\begin{hbtheorem}{thm:p15-w02b-p15-025-coverage}{proved=annex:p15-w02b-p15-025-coverage}{What the confidence coefficient is, and is not, a confirmation of}
In the setting of \cref{def:p15-w02b-p15-025-joint}:
\begin{enumerate}
\item (Direct inference.) If \(\pi(\theta_0)>0\), then
\(Q_\pi(\{(\theta,x):\theta\in C(x)\}\mid\theta=\theta_0)=\gamma_C(\theta_0)\).
\item (Average.) For every \(\pi\),
\[
 \sum_{x:\,Q_\pi(x)>0}Q_\pi(x)\,Q_\pi\bigl(C(x)\bigm|x\bigr)
 =\sum_{\theta}\pi(\theta)\gamma_C(\theta).
\]
Hence if \(\gamma_C\ge1-\alpha\) on \(\Theta\), the
\(Q_\pi\)-expected confirmation of the reported set relative to the
observed sample is at least \(1-\alpha\), with equality for every
\(\pi\) when the coverage is exactly \(1-\alpha\).
\item (No sample-wise reading.) Let \(\Theta=\{0,1\}\),
\(\mathsf X=\{1,2,3\}\),
\(P_0=(\frac9{10},\frac1{10},0)\), \(P_1=(0,\frac1{10},\frac9{10})\), and
\(C(1)=\{0\}\), \(C(2)=\varnothing\), \(C(3)=\{1\}\). The coverage is
exactly \(9/10\) at both parameter values. For every \(\pi\) with
\(\pi(0),\pi(1)>0\), the confirmation of the reported set relative to
the observed sample is \(1\) at \(x=1\) and \(x=3\) and \(0\) at
\(x=2\); it is never \(9/10\).
\end{enumerate}
\end{hbtheorem}

Here \(Q_\pi(C(x)\mid x)\) abbreviates
\(Q_\pi(\{\theta\in C(x)\}\times\{x\})/Q_\pi(x)\). Item~1 is Carnap's
statement that the confidence coefficient equals the confirmation of
the interval statement relative to a parameter hypothesis. Item~3 is an
exact model of his warning that it cannot be read as the confirmation
relative to the observed sample
\cite[sec.~98, p.~517]{Carnap1950Probability}; no prior repairs it,
since the content at \(x=2\) is zero for all of them. Item~2 states the
one relation that does hold: the coefficient bounds an average of
confirmations over samples, taken under the confirmation function's own
predictive distribution. The interval case with a continuous parameter
and the uniform prior is in \cref{thm:p15-confidence-credibility-boundary}.

\section{Wald's minimax rule and weighted loss}
\label{sec:p15-w02b-p15-025-minimax}

\begin{hbtheorem}{thm:p15-w02b-p15-025-minimax}{proved=annex:p15-w02b-p15-025-minimax}{The two rules disagree on an interval of confirmation values}
Let the parameter take values \(\theta_1,\theta_2\), let the available
actions be \(a_1,a_2\) and their randomizations \(\delta_t\), which take
\(a_1\) with probability \(t\) and \(a_2\) otherwise
(\(0\le t\le1\)), and let the losses be
\[
 L(a_1,\theta_1)=0,\quad L(a_1,\theta_2)=10,\quad
 L(a_2,\theta_1)=L(a_2,\theta_2)=4 .
\]
\begin{enumerate}
\item The unique minimax choice is \(\delta_0=a_2\), with maximum
expected loss \(4\).
\item If the evidence gives \(\theta_2\) the confirmation \(q\), the
confirmation-weighted loss of \(\delta_t\) is \(4+t(10q-4)\). It is
minimized uniquely by \(a_1\) when \(q<2/5\), uniquely by \(a_2\) when
\(q>2/5\), and by every \(\delta_t\) when \(q=2/5\).
\end{enumerate}
So the two rules choose differently for every \(q\in[0,2/5)\), and
agree for \(q\in(2/5,1]\).
\end{hbtheorem}

No sample is needed for the disagreement, which already arises in the
decision made after the data have fixed \(q\). This is an exact case of
Carnap's remark that Wald's rule leads in general to decisions
different from those of his own rules
\cite[sec.~98, p.~517]{Carnap1950Probability}. The minimax choice
ignores \(q\) altogether; the weighted rule uses \(q\) and the losses
and nothing else.

\begin{hbwarning}{def:p15-w02b-p15-025-stat-limits}{What the statistical comparisons do not settle}
\begin{enumerate}
\item Each comparison is conditional on a confirmation function, or on
values \(\pi\) for the parameter hypotheses. Carnap's \(c^*\) is defined
only for the simple languages of his book, and he leaves its extension
to quantitative languages open
\cite[sec.~98, p.~519]{Carnap1950Probability}. Outside those languages,
the confirmation side of every result above uses an assumed prior, not
one derived from \(c^*\).
\item The comparisons of outputs that Carnap promises for his second
volume are not reconstructed here
\cite[sec.~98, pp.~515, 518--520]{Carnap1950Probability}. The results of
this chapter show that the outputs differ, and by how much, in the
stated models; they do not show which output is better.
\item The descriptions of Fisher, Neyman, Pearson and Wald in this
chapter are Carnap's 1950 descriptions. Their own constructions are
presented in
Chapters~\ref{sec:p15-fisher-estimation},
\ref{sec:p15-hypothesis-tests-confidence}
and~\ref{sec:p15-neyman-confidence-constructions}, and Carnap's account
is not a description of later statistics.
\item No empirical claim is made about how statisticians estimate or
decide.
\end{enumerate}
\end{hbwarning}

\section{Relationships and sources}
\label{sec:p15-w02b-p15-025-relations}

The constructions of systematic power used here are those of
Chapter~\ref{sec:p15-systematic-power}, in particular
\cref{thm:p15-power-optimal-split,thm:p15-power-probability-dual,thm:p15-power-dualism}.
The relevance relations and measures are those of
Chapter~\ref{sec:p15-rationality-4-3}. Carnap's covariance measure
\(\Delta_P\), which reappears in
\cref{thm:p15-w02b-p15-025-increment} as the numerator of
\(s_r-g\), is compared with Merin's measures in
\cref{sec:p15-w02b-p15-024-carnap}. Hempel's prediction criterion and its
defects are in Chapter~\ref{sec:p15-hempel-confirmation-and-paradoxes};
the deductive-nomological pattern to which Part~IV is attached is in
\cref{sec:p16-dn-question}. Merin attributes the bounded likelihood
transform KO to Kemeny and Oppenheim in 1952
\cite[p.~270]{Merin2021Relevance}; that paper is not presented here.

For the statistical comparison, the \(c\)-mean estimate, its
reliability and Carnap's general estimation programme are in
\cref{sec:p15-carnap-general-estimation}; the predictive value \(c^*\)
is \cref{thm:p15-carnap-predictive}; Fisher's binomial and regularity
results are in Chapter~\ref{sec:p15-fisher-estimation}; tests,
confidence procedures and the coverage--posterior boundary are in
Chapter~\ref{sec:p15-hypothesis-tests-confidence}; and the normal model
with unknown mean and variance is in \cref{def:p15-student-model}.

The sources are Hempel and Oppenheim, Part~IV, Sections~8 and~9
\cite[pp.~164--173]{HempelOppenheim1948Explanation}; Merin's
chapter~4.3, Sections~1--5
\cite[pp.~265--273]{Merin2021Relevance}; and Carnap's Section~98
\cite[sec.~98, pp.~511--520]{Carnap1950Probability}. The identities,
equivalences and models of this chapter are not stated in these
sources; they are proved in the annex.

\index{systematic power!and evidential relevance}
\index{evidential relevance!systematic power}
\index{prediction criterion!systematic power}
\index{old evidence!logical range measure}
\index{maximum likelihood!and predictive confirmation}
\index{confidence coefficient!and degree of confirmation}
\index{minimax rule!and confirmation-weighted loss}
